% Curriculum Vitae -- Ricky Ding
% Compile with pdfLaTeX (e.g. on Overleaf).

\documentclass[11pt,a4paper]{article}

\usepackage[T1]{fontenc}
\usepackage[utf8]{inputenc}
\usepackage{lmodern}
\usepackage[margin=2.2cm]{geometry}
\usepackage{enumitem}
\usepackage{titlesec}
\usepackage{xcolor}
\usepackage{hyperref}

\hypersetup{colorlinks=true, urlcolor=blue, linkcolor=blue, pdfauthor={Ricky (Ruiqi) Ding}, pdftitle={Curriculum Vitae}}

% Supporting documents hosted on the website
\newcommand{\doc}[2]{\href{https://neuroai-research.github.io/Ricky/#1}{#2}}

\pagestyle{empty}
\setlength{\parindent}{0pt}
\setlist[itemize]{leftmargin=1.2em, itemsep=1pt, topsep=2pt, parsep=0pt}

\titleformat{\section}{\large\scshape}{}{0em}{}[\vspace{-0.6ex}\rule{\textwidth}{0.4pt}]
\titlespacing*{\section}{0pt}{1.6ex}{1.0ex}

% \entry{Title}{Dates}{Subtitle}{Location}
\newcommand{\entry}[4]{%
  \textbf{#1} \hfill #2\\
  \textit{#3} \hfill \textit{#4}\par\vspace{2pt}}

\begin{document}

% ------------------------------------------------------------------
\begin{center}
  {\LARGE \textsc{Ricky (Ruiqi) Ding}}\\[4pt]
  \href{https://neuroai-research.github.io/Ricky/}{neuroai-research.github.io/Ricky}
  \;$\cdot$\; \href{https://github.com/tesla-cat}{github.com/tesla-cat}\\
  \href{mailto:e0134117@u.nus.edu}{e0134117@u.nus.edu}
  \;$\cdot$\; \href{https://www.linkedin.com/in/rickyding1997}{linkedin.com/in/rickyding1997}
\end{center}

% ------------------------------------------------------------------
\section{Research Interests}

I am broadly interested in AI and neuroscience/neurotechnology research. My
current self-directed focus is neuroscience-inspired AI: understanding how
specific modules and circuits of the human brain achieve their functions, and
replicating those functions with modern AI methods, in the spirit of
\href{https://doi.org/10.1016/j.neuron.2017.06.011}{Hassabis et al.\ (\textit{Neuron}, 2017)}.
I trained in physics and have since worked as a software developer.

% ------------------------------------------------------------------
\section{Current Research Project}

\textbf{NeuroAI Research} \hfill 2026 -- present\\
(\href{https://neuroai-research.github.io/}{website},
\href{https://github.com/NeuroAI-Research/neu-ai}{code})
\begin{itemize}
  \item Self-directed study towards replicating brain modules and circuits
        with modern AI methods; currently in the learning and
        re-implementation stage
  \item Minimal re-implementations of models from Dayan \& Abbott,
        \textit{Theoretical Neuroscience}: neural encoding and decoding,
        information theory, network models, plasticity, reinforcement learning
  \item Minimal re-implementations of GPT-2, Gemma~2, PPO, SAC, DreamerV3
        and flow matching
\end{itemize}

% ------------------------------------------------------------------
\section{Education}

\entry{National University of Singapore}{2016 -- 2020}
      {B.Sc.\ (Honours) in Physics, Highest Distinction}{Singapore}
\begin{itemize}
  \item Thesis (supervisor: Prof.\ Berthold-Georg Englert):
        \doc{2020-03-01_BSc_Thesis.pdf}{\textit{Density Potential Functional
        Theory in $x$ and $p$ spaces and its implementation}}
  \item Scholarship: \href{https://nus.edu.sg/oam/scholarships/scholarships-for-freshmen-(international-students)/science-technology-undergraduate-scholarship}{Science \& Technology Undergraduate Scholarship}
  \item Degree certificate: \doc{2020-06-30_NUS.png}{image};
        \doc{2020-06-30_NUS.opencert}{OpenCert}, verifiable at
        \href{https://www.opencerts.io}{opencerts.io}
\end{itemize}
\vspace{4pt}

\entry{University of G\"ottingen}{2019}
      {Undergraduate exchange programme (non-degree)}{G\"ottingen, Germany}
\begin{itemize}
  \item Took courses in biophysics from the M.Sc.\ Physics programme
        (\doc{2019-08-06_Gottingen.pdf}{certificate})
\end{itemize}

% ------------------------------------------------------------------
\section{Research Experience}

\entry{Centre for Quantum Technologies, National University of Singapore}{2020 -- 2021}
      {Research Assistant (Software Development)}{Singapore}
\begin{itemize}
  \item Compared deep reinforcement learning with gradient ascent pulse
        engineering (GRAPE) for quantum optimal control
        (\doc{2020_09_01_RL_Quantum_Optimal_Control.pdf}{report})
  \item Studied transmon (superconducting) quantum computing
        (\doc{2020-06-22_Transmon_Quantum_Computer.pdf}{notes})
  \item Worked on IQ mixer tuning
        (\doc{2021_01_01_IQ_Mixer_Tuning.pdf}{notes})
  \item Employment \doc{2021-05-19_NUS_CQT_Certificate.pdf}{certificate}
\end{itemize}
\vspace{4pt}

\entry{Research Internships, National University of Singapore}{2018 -- 2019}
      {Undergraduate research}{Singapore}
\begin{itemize}
  \item 2019: \href{https://ions-sg.org/index.html}{Quantum Computing Lab},
        supervised by
        \href{https://scholar.google.com/citations?user=B4X3nCMAAAAJ&hl=en}{Prof.\ Dzmitry Matsukevich}
  \item 2018: Quantum Theory Group, supervised by
        \href{https://scholar.google.com/citations?user=ttUoRO4AAAAJ&hl=en}{Prof.\ Berthold-Georg Englert}
\end{itemize}

% ------------------------------------------------------------------
\section{Professional Experience}

\entry{Frontline Industrial Software}{2022 -- 2025}
      {Software Developer}{Singapore}
\begin{itemize}
  \item Developed software products for industrial clients, applying:
  \begin{itemize}
    \item large language models for content generation
    \item vector embeddings and similarity search
    \item automatic differentiation and stochastic gradient descent (Adam)
    \item density-based clustering (\href{https://github.com/scikit-learn-contrib/hdbscan}{HDBSCAN})
    \item mixed-integer linear programming (\href{https://github.com/coin-or/pulp}{PuLP})
  \end{itemize}
  \item Frontend: Vue, React, Capacitor (Web, Android, iOS);
        backend: Python, JavaScript (FastAPI, Bun)
\end{itemize}

% ------------------------------------------------------------------
\section{Personal Projects}

\textbf{Algorithmic Trading}
\begin{itemize}
  \item Differentiable backtesting; deep reinforcement learning
\end{itemize}
\vspace{4pt}

\textbf{Tesla Coil Gun} (\href{https://www.youtube.com/watch?v=l_bg02wG8IY}{video})
\begin{itemize}
  \item \href{https://github.com/TeslaCoilResearch/tesla-coil-simulator}{Physics simulation}
        of a dual-resonant solid-state Tesla coil (DRSSTC);
        \doc{TC_Music_App.png}{waveform synthesiser}
  \item \doc{2023-12-16_TC_Gun_driver.pdf}{Schematic} and
        \doc{2023-12-16_driver_PCB_3D.png}{PCB} design;
        CPLD programming in Verilog
\end{itemize}

\end{document}
